\documentclass{article}
\usepackage[T1]{fontenc}
\usepackage{lmodern}
\usepackage{graphicx}
\usepackage{amsmath,amssymb}
\usepackage[a4paper,margin=2cm]{geometry}
\usepackage[absolute,overlay]{textpos}
\setlength{\TPHorizModule}{1mm}
\setlength{\TPVertModule}{1mm}
\usepackage[indonesian]{babel}
\usepackage{hyperref}
\setlength{\emergencystretch}{2em}
% unit: Halaman 23

\begin{document}

% === Halaman 23 (llm) ===

\section*{Kelemahan RSA}

\begin{itemize}
    \item RSA lebih lambat daripada algoritma kriptografi kunci-simetri seperti $DES$ dan $AES$
    \item Dalam praktek, $RSA$ tidak digunakan untuk mengenkripsi pesan yang berukuran besar, tetapi digunakan untuk mengenkripsi kunci simetri (kunci sesi) dengan menggunakan kunci publik penerima pesan. Karena, kunci simeteri umumnya relative pendek.
    \item Pesan tetap dienkripsi dengan algoritma simetri seperti $DES$ atau $AES$.
    \item Pesan dan kunci simetri (yang terenkripsi) dapat dikirim bersamaan. Penerima mendekripsi kunci simetri dengan kunci privatnya, lalu mendekripsi pesan dengan kunci simetri tersebut.
    \item Kombinasi kedua system kriptografi ini (simeteri dan nirsimetri) dinamakan \textit{hybrid cryptography}.
\end{itemize}

\end{document}
